\section*{Заключение}

В ходе лабораторной работы реализована и исследована модель арифметического конвейера для попарного умножения компонентов двух векторов чисел (вариант 2: умножение 4-разрядных чисел с младших разрядов со сдвигом частичной суммы вправо). Разрядность операндов $p = 4$, число стадий конвейера $n = 4$, ранг задачи $r$ совпадает с количеством пар $m$.

Реализованы ввод векторов (генерация и ручной ввод), проверка разрядности, пошаговый вывод состояния конвейера по тактам с отображением данных в двоичном виде и времени, а также расчёт результирующего вектора. Для сбалансированного конвейера получены формулы времени выполнения $T_1 = r n$, $T_n = n + r - 1$, коэффициента ускорения $K_y(n, r) = r n / (n + r - 1)$ и эффективности $e(n, r) = r / (n + r - 1)$. Построены семейства графиков зависимостей $K_y$ и $e$ от ранга $r$ и от числа стадий $n$; объяснены асимптоты и характер изменения графиков при изменении параметров.

Установлено соотношение между рангами при одинаковой эффективности: при $e(n_1, r_1) = e(n_2, r_2)$ и $n_1 > n_2$ выполняется $r_1 > r_2$. Получена формула порогового ранга $r_0$ для выполнения условия $e(n, r_0) > e_0$: $r_0 > e_0(n - 1) / (1 - e_0)$. Рассмотрены способы перестройки несбалансированного конвейера для достижения заданной эффективности и максимальной скорости; показано, что максимальная скорость достигается балансировкой конвейера (выравниванием времени стадий), при этом для полученного конвейера справедливы приведённые формулы $K_y(n, r)$ и $e(n, r)$.

Цель работы достигнута: модель конвейера реализована, протокол работы получен, теоретические зависимости проверены и оформлены в отчёте.
